\section{Problem formulation}\label{sec:problem}

\subsection{Problem statement}\label{sec:setting}

\sumpara The line, the jobs, and the planner's commitments define the problem. \sumparaend
\fin A line of $m$ stages, indexed by $k$ in line order, processes $n$ jobs, indexed by $j$, and each stage has one machine. \fine
\fin Job $j$ is released at time $r_j$, is due at $d_j$, carries a weight $w_j$ that scores its priority, and requires a processing time $p_{jk} > 0$ at stage $k$. \fine
\fin Every job visits every stage in line order without interruption, jobs cannot pass one another, so one sequence of the jobs holds at every stage, and the buffers between stages are unlimited. \fine
\fin We assume that every stage is idle at time zero, that the planner fixes the sequence before the first job starts and does not revise it, and that among sequences with equal weighted tardiness the planner prefers less waiting between stages under earliest-start execution. \fine

\subsection{Earliest-start recursion}\label{sec:timing}

\sumpara A sequence and the earliest-start recursion determine the schedule and both measures. \sumparaend
\fin A sequence $\pi$ lists the jobs in the order in which they enter the line, with $\pi(i)$ the job in position $i$. \fine
\fin Given $\pi$, each operation starts at the later of the job's completion at the previous stage and the stage's completion of the preceding job: \fine
\begin{equation}\label{eq:recursion}
S_{\pi(i),k} = \max\bigl\{C_{\pi(i),k-1},\; C_{\pi(i-1),k}\bigr\}, \qquad C_{\pi(i),k} = S_{\pi(i),k} + p_{\pi(i),k},
\end{equation}
\fin with the conventions $C_{\pi(i),0} = r_{\pi(i)}$ and $C_{\pi(0),k} = -\infty$. \fine
\fin The tardiness of job $j$ is $T_j = \max\{0,\, C_{jm} - d_j\}$, and the primary and secondary measures are the total weighted tardiness and the total waiting between stages, \fine
\begin{equation}\label{eq:measures}
Z(\pi) = \sum_{j=1}^{n} w_j T_j, \qquad W(\pi) = \sum_{j=1}^{n} \sum_{k=2}^{m} \bigl(S_{jk} - C_{j,k-1}\bigr),
\end{equation}
\fin in weighted minutes and job-minutes, respectively. \fine
\fin Both measures are evaluated on the earliest-start schedule of \eqref{eq:recursion}, under which a job waits between stages only while the next stage is still occupied by the preceding job. \fine

\sumpara The recursion yields the smallest tardiness of any schedule with that sequence, so the problem reduces to choosing a sequence. \sumparaend
\fin By induction over positions and stages, any feasible schedule that keeps sequence $\pi$ starts each operation no earlier than \eqref{eq:recursion}, so the recursion gives every job its earliest completion at the last stage and $\pi$ its smallest value of $Z$. \fine
\fin The planning problem is therefore to choose the sequence that minimizes $Z(\pi)$, which is the permutation flow shop problem with release dates and total weighted tardiness. \fine
\fin Applied with realized processing times, \eqref{eq:recursion} is also the right-shift execution of a fixed sequence described by \citet{aytug2005executing}, which we use in the stress test. \fine

\subsection{Positional mixed-integer formulation}\label{sec:mip}

\sumpara A positional program states the same problem in a form that a solver can certify. \sumparaend
\fin To certify optimal values, we state the problem as a mixed-integer program over sequence positions $q = 1, \ldots, n$. \fine
\fin The binary variable $x_{jq}$ equals 1 when job $j$ occupies position $q$, $B_{qk}$ and $E_{qk}$ are the start and end of the job in position $q$ at stage $k$, and $T_j \geq 0$ is the tardiness of job $j$: \fine
\begin{align}
\min \quad & \sum_{j=1}^{n} w_j T_j \label{eq:mip-obj}\\
\text{s.t.} \quad & \textstyle\sum_{q=1}^{n} x_{jq} = 1 \quad \forall j, \qquad \sum_{j=1}^{n} x_{jq} = 1 \quad \forall q, \label{eq:mip-assign}\\
& E_{qk} = B_{qk} + \textstyle\sum_{j=1}^{n} p_{jk}\, x_{jq} \quad \forall q, k, \label{eq:mip-proc}\\
& B_{q1} \geq \textstyle\sum_{j=1}^{n} r_j\, x_{jq} \quad \forall q, \label{eq:mip-release}\\
& B_{qk} \geq E_{q,k-1} \quad \forall q,\ k \geq 2, \label{eq:mip-stage}\\
& B_{qk} \geq E_{q-1,k} \quad \forall q \geq 2,\ k, \label{eq:mip-machine}\\
& T_j \geq E_{qm} - d_j - U_j\,(1 - x_{jq}) \quad \forall j, q, \label{eq:mip-tardy}\\
& x_{jq} \in \{0,1\}, \quad T_j \geq 0. \notag
\end{align}
\fin Constraints \eqref{eq:mip-assign} place every job in exactly one position, \eqref{eq:mip-proc} processes each operation without interruption, \eqref{eq:mip-release} and \eqref{eq:mip-stage} enforce release dates and stage order, and \eqref{eq:mip-machine} allows a stage to start a job only after it finishes the preceding position, which also keeps one order at every stage. \fine
\fin Constraint \eqref{eq:mip-tardy} assigns the tardiness of the job in position $q$ to that job and is inactive for every other job when $U_j = \max\{0,\, \max_h r_h + \sum_{h,k} p_{hk} - d_j\}$, where $h$ ranges over all jobs. \fine

\sumpara The tardiness bound never excludes the optimal schedule, so the program's optimum equals the best sequence's tardiness. \sumparaend
\fin Under \eqref{eq:recursion}, every job of every sequence finishes by $\max_h r_h + \sum_{h,k} p_{hk}$, because each completion time is a release date plus the processing times along one path through the schedule. \fine
\fin The bound $U_j$ therefore never forces a positive tardiness on the earliest-start schedule of any sequence, which by \cref{sec:timing} is the best schedule for that sequence, and the optimal value of the program equals $\min_\pi Z(\pi)$. \fine
\fin The program serves only to certify the optima of the 6-job and 15-job instances and is never one of the compared methods. \fine
